We diagnose the LCO FLOYDS pipeline \banzai{} for reverberation mapping (RM), using a Fairall~9 campaign at Siding
Spring (COJ; 100 spectra through the 6\,arcsec slit, 57 standard-star spectra, 122 arcs) and an NGC\,4151 campaign at
Haleakal\=a (OGG; 60 spectra, 100 standard-star spectra), and we describe target-independent corrections, now built into
\banzai{} and run from the raw frames. Our final version (improved v3) addresses every problem we found.
(i) \textit{Wavelengths}: the arc fit leaves H$\beta$/[O\,{\sc iii}] and H$\alpha$ unconstrained (the archive red order
is off by $-8$\,\AA{} at Hg\,4916) and a compact object's zero point in a wide slit depends on its position across the
slit (up to $\pm2''$, i.e.\ $\pm17$\,\AA{} for NGC\,4151). Adding the detected but unused lines to the arc fit and
measuring the object's zero point from the telluric bands in its own light bring the [O\,{\sc iii}] scatter from
exposure to exposure to 0.5--0.9\,\AA{} and the two orders into agreement to 0.3\,\AA{} (COJ). (ii)
\textit{Fringing}: the archive frames are divided by an eight-month-old lamp flat shifted along the slit by up to
hundreds of pixels, which destroys data in 13 of 50 frames and biases the flux by up to $\pm10$ per cent; a lamp flat
gives the phase of the fringes but not their shape. A two-harmonic fringe model with the lamp phase, fitted to each
extracted spectrum jointly with a telluric template, reduces the red-order fringing of Fairall~9 from 13.7 (archive)
and 8.1 (\banzai{}~1.1.0) to 1.9 per cent. (iii) \textit{Telluric and flux calibration}: standards are selected by
proposal rather than slit width, several nightly standards are missing from the catalogue, a single standard
transfers its clouds and its noise to the science, and the archive spectra are twice too bright because 1.0.1 does not
scale by the exposure time. An ensemble of the nine nearest standards with the grey level of the clearest half, a
telluric correction confined to the absorption bands and two added standards reduce the scatter of the red/blue
[O\,{\sc iii}] flux ratio from 29.2 (archive) and 8.9 (1.1.0) to 6.9 per cent, and bridge a factor-1.8 throughput change
of the OGG telescope. (iv) \textit{Traces}: a physical model of the object trace (refraction plus an instrumental
shape for each site) keeps the faint ends of the orders at 0.1--0.3\,px where the polynomial fails. The uncertainties
are correct per pixel at both sites. We list the open issues and give recommendations for observing and for the
pipeline.
